Les équations \eqref{eq:navier_stokes} peuvent être simplifiées sous les
hypothèses suivantes :
\begin{itemize}
    \item l'écoulement est incompressible :
    $\nabla \cdot \mathbf{u} = 0$,

    \item le fluide est homogène et sa masse volumique est constante :
    $\rho(\mathbf{x}, t) = \rho_0 > 0$,

    \item la viscosité dynamique est constante : $\mu = \mu_0$,

    \item la gravité est négligée : $\mathbf{g} = \mathbf{0}$.
\end{itemize}
Pour un fluide newtonien à viscosité constante, la contrainte
d'incompressibilité donne
$\nabla \cdot \boldsymbol{\tau} = \mu_0 \nabla^2 \mathbf{u}$.
En posant $\nu = \mu_0 / \rho_0$, la viscosité cinématique,
on obtient les équations de Navier--Stokes incompressibles :
\begin{equation}
    \label{eq:navier_stokes_incompressible}
    \boxed{
    \begin{aligned}
        \nabla \cdot \mathbf{u} &= 0, \\
        \frac{\partial \mathbf{u}}{\partial t}
        + \nabla \cdot \left(\mathbf{u} \otimes \mathbf{u}\right)
        &= -\frac{1}{\rho_0}\nabla p + \nu \nabla^2 \mathbf{u}.
    \end{aligned}}
\end{equation}

L'écoulement considéré est bidimensionnel :
la vitesse s'écrit $\mathbf{u} = (u,v)$ et les champs
dépendent de $(x,y,t)$, sans dépendance selon $z$.

